% =====================================================================
% Capítulo 3 — Esboço de seção/interlúdio (análise de rede textual)
% Modelo: esboco_secao_rede_textual.tex (capítulo 1).
%
% Estilo seguido (regras da tese): primeira pessoa do singular, sem
% travessões em texto principal, "capítulo" minúsculo em referências
% cruzadas, aspas com \enquote{}, sem fórmulas do tipo "não X, mas Y".
%
% Observação técnica: copie
%   rede_textual/cap3/infranodus_cap3_focus.png
% para
%   figuras/cap.3/infranodus_cap3_focus.png
% =====================================================================

\subsection*{Interlúdio: a rede textual do capítulo}

Aplico de novo a análise de rede textual introduzida no
capítulo~\ref{capitulo1}, agora sobre os 18.887~\textit{tokens}
deste capítulo. O grafo de 180 termos separa o arranjo do C4AI em três
planos que reconheço como as três camadas do argumento. Há a
infraestrutura histórica das máquinas de calcular, reunida em
\textit{tecnologia}, \textit{hollerith}, \textit{tabulação},
\textit{máquina} e \textit{trajetória}, a longa duração que vai das
máquinas de cartão perfurado até a IBM. Há o arranjo público-privado do
presente, em \textit{público}, \textit{arranjo}, \textit{corporação},
\textit{universidade} e \textit{brasil}. E há o vocabulário do método,
que forma comunidade própria em torno de \textit{rede}, \textit{seguir},
\textit{ator}, \textit{associação} e \textit{descrevo}: o gesto de
\enquote{seguir os atores} tem, no grafo, existência lexical autônoma,
distinta tanto da instituição descrita quanto da história que a
antecede.

Os termos de maior \textit{degree} são \textit{cláudio},
\textit{pesquisa}, \textit{rede}, \textit{fábio} e \textit{centro}, e as
maiores pontes por \textit{betweenness} são \textit{pesquisa}
(0,31), \textit{centro} (0,26) e \textit{cláudio}
(0,22), seguidas de \textit{rede} e \textit{fábio}. Os dois
interlocutores que sigo no campo aparecem como uma dupla de forte
associação, \textit{cláudio}~$\leftrightarrow$~\textit{fábio}, e ocupam
posição de intermediação: são, na rede, os nós por onde a etnografia
passa. \textit{Centro} e \textit{pesquisa} funcionam como termos que
circulam entre a história infraestrutural, o arranjo institucional e o
trabalho de campo, costurando os três planos. O objeto técnico do centro
também deixa rastro nítido, nos pares
\textit{linguagem}~$\leftrightarrow$~\textit{natural} (0,75),
\textit{processamento}~$\leftrightarrow$~\textit{natural} (0,70),
\textit{código}~$\leftrightarrow$~\textit{aberto} (0,73) e
\textit{línguas}~$\leftrightarrow$~\textit{indígenas} (0,78), que
condensam as frentes de pesquisa que descrevo sem que eu precise
enumerá-las.

\begin{figure}[!htb]
\centering
\includegraphics[width=1.0\textwidth,keepaspectratio]{figuras/cap.3/infranodus_cap3_focus.png}
\caption[Rede textual do capítulo~\ref{capitulo3} (\textit{text network
analysis})]{Rede textual do capítulo~\ref{capitulo3}
(\textit{text network analysis}). Os nós são os termos de maior
\textit{degree} ponderado; o tamanho do nó acompanha o \textit{degree};
a cor indica a comunidade detectada por Louvain ponderado; as arestas,
arqueadas, recebem a mistura das cores das comunidades que ligam.
Versão de alta resolução, arquivo \texttt{.gexf} para o \textit{Gephi} e
relatório completo em \texttt{rede_textual/cap3/} no repositório desta
tese.}
\label{fig:infranodus-cap3}
\end{figure}

A leitura das ligações fracas devolve ao capítulo a sua própria
promessa. A costura mais rala está entre o vocabulário do centro de
pesquisa (\textit{pesquisa}, \textit{centro}, \textit{científico},
\textit{fapesp}) e o vocabulário do método ator-rede (\textit{rede},
\textit{ator}, \textit{associação}): a instituição que descrevo e o
método com que a descrevo caminham lado a lado, ainda pouco entrelaçados
na letra do texto. Fraca também é a ligação entre a infraestrutura
histórica da IBM e da tabulação Hollerith e o centro do presente, o que
me parece a junta decisiva a fortalecer, já que o argumento do capítulo
sustenta que a racionalidade sedimentada nas máquinas de calcular
continua a operar no arranjo contemporâneo. Tornar essa continuidade
lexicalmente explícita, ligando \textit{hollerith} e \textit{tabulação}
a \textit{centro} e \textit{arranjo}, é o trabalho que a rede me aponta.
A leitura latouriana ampara esse diagnóstico, pois seguir a associação
\parencite{latour2005reassembling} é justamente reconstruir os elos que a
descrição corrente deixa implícitos, e a rede me mostra onde eles ainda
faltam.
